
@@M0@@
où @@M1@@, @@M2@@
sont les fibres \emph{géométriques} de @@M3@@. D'ailleurs, on sait que maintenant l'homomorphisme @@M4@@ de spécialisation pour la cohomologie est un isomorphisme (SGA 4 XVI 2.2), donc la commutativité du carré précédent peut s'exprimer encore par celle du carré qui s'en déduit en inversant la flèche précédente. Passant alors à la limite projective sur @@M5@@ dans les carrés commutatifs ainsi obtenus, on trouve un carré commutatif
@@M6@@

\textbf{7.13.5.} Une autre façon d'exprimer cette commutativité est de dire que, si @@M7@@ est un élément de @@M8@@ et @@M9@@ son image dans @@M10@@, alors @@M11@@ et @@M12@@ sont les valeurs, en les points géométriques @@M13@@ et @@M14@@ respectivement, d'une \emph{section} de @@M15@@. Cette section s'interprète d'ailleurs aisément, en prenant un élément @@M16@@ qui induit @@M17@@, d'où un @@M18@@, qui définit la section cherchée de @@M19@@.

\textbf{7.13.6.} Soit @@M20@@ un schéma. Alors par passage aux gradués associés, l'homomorphisme (7.13.0) définit des homomorphismes
@@M21@@
l'image par @@M22@@ d'un élément @@M23@@ de @@M24@@ étant par définition @@M25@@. (Le fait que (7.13.0) est compatible avec la @@M26@@-filtration sur 
